% ================= TÓM TẮT =================
\newpage
\thispagestyle{plain}

\begin{center}
    {\bfseries\Large TÓM TẮT}
\end{center}

\vspace{1cm}

Chuyên đề giải quyết hai bài toán xử lý chuỗi có cấu trúc khác nhau: tìm mọi vị trí xuất hiện
của một mẫu trong văn bản và gợi ý các từ có chung tiền tố trong một từ điển. Bài toán thứ nhất
được giải bằng thuật toán Knuth--Morris--Pratt; bài toán thứ hai sử dụng cấu trúc dữ liệu Trie.

Phần cài đặt gồm hai mô-đun độc lập, tương thích Python 3.10+ và chỉ dùng thư viện chuẩn.
Mô-đun KMP công khai prefix function và trả về mọi vị trí khớp, kể cả các lần xuất hiện chồng
lấn. Mô-đun Trie hỗ trợ chèn, tìm chính xác, truy vấn tiền tố và autocomplete có giới hạn số
kết quả. Phần mở rộng cài đặt thêm thuật toán Aho--Corasick cho bài toán nhiều mẫu.

Tính đúng đắn được lập luận bằng quy nạp và bất biến vòng lặp. Bộ kiểm thử gồm 276 test tự động
và 10 doctest, có baseline độc lập, test biên, test đối kháng và đối chiếu ngẫu nhiên với seed
cố định. Năm thí nghiệm được chạy trên dữ liệu từ \(10^3\) đến \(10^6\) ký tự và tới
\(5\cdot10^5\) từ.

Kết quả cho thấy số phép so sánh của KMP trên đầu vào đối kháng đạt tối đa \(1{,}9999n\), sát
chặn \(2n\). Với danh sách đã sắp gồm 500.000 từ, Trie nhanh hơn baseline quét tuyến tính
1.062 lần trong đúng cấu hình đo. Exact search chỉ duyệt từ 6 đến 10 nút, bằng độ dài từ truy
vấn. Kết quả bộ nhớ đồng thời cho thấy lợi ích của Trie phụ thuộc mức chia sẻ tiền tố và chi
tiết cài đặt.

\vspace{0.5cm}
\noindent\textbf{Từ khoá:} KMP, prefix function, Trie, autocomplete, Aho--Corasick,
phân tích khấu hao, kiểm thử đối kháng.
